\chapter{System Analysis and Requirements}

\section{Requirement Analysis}

Nexis is a role-based, multi-tenant platform. Every piece of data in the
system belongs to exactly one organisation (tenant), and every request is
evaluated against both the caller's organisation and the caller's role within
it. The requirement analysis therefore begins by identifying the actors who
interact with the system boundary.

Three \emph{human} actor types were identified:

\begin{itemize}[nosep]
  \item \textbf{Owner/Admin} --- the organisation's administrator, responsible
        for membership, roles, billing, policy, integrations, and the final
        human decision on high-stakes automated repairs.
  \item \textbf{Engineer/Member} --- the working engineer who investigates
        incidents, exercises the evaluation tooling, and interacts with the
        recovery pipeline day to day.
  \item \textbf{Viewer} --- a read-only observer (for example a stakeholder or
        an auditor) who may inspect the system's state but may not change it.
\end{itemize}

The role model is \emph{additive}: an Admin holds every permission an Engineer
holds, and an Engineer holds every permission a Viewer holds. Consequently,
the functional requirements in Section~3.2 list only the permissions each role
\emph{adds} to the tier below it, without repetition. (Internally the database
role enumeration is owner/admin/member; the console additionally presents the
read-only Viewer experience through RBAC gating on the frontend and API.)

Unusually for a requirement analysis, a fourth actor category must be modelled
that is not human at all: the \textbf{nine-agent self-healing fleet}. Most
requirement-analysis chapters enumerate only human actors, because in a
conventional system every state change ultimately traces back to a person
clicking something. That assumption does not hold for Nexis. The agent fleet
acts autonomously inside the system boundary: it detects faults without being
asked, reads and writes the same PostgreSQL, Neo4j, and object-storage data
stores the human-facing console uses, and triggers real external side effects
--- opening GitHub pull requests, applying rollbacks through Argo
CD~\cite{argocd}, and redeploying preview containers. An actor model that
omitted the fleet would misrepresent the system: roughly half of all writes to
the incident, patch, and audit tables originate from agents rather than
people. Treating the fleet as a first-class actor also makes the safety
analysis explicit --- every autonomous capability listed in Section~3.2.4 is
bounded by a corresponding control (sandboxing, severity classification, or
the human approval gate) listed in Section~3.3.

The Level~0 and Level~1 use case diagrams for these four actors are presented
together with the rest of the system's UML design in
Section~\ref{sec:usecase-diagrams} (Figures~\ref{fig:usecase-l0} and
\ref{fig:usecase-l1}), alongside the data flow diagrams that show how the same
actors move data through the platform.

\section{Functional Requirements}

The functional requirements are grouped by actor. Because the role model is
additive, each human tier is listed as a delta over the tier below it.

\subsection{Owner/Admin can:}

\begin{itemize}[nosep]
  \item Do everything an Engineer/Member and a Viewer can do (role
        additivity).
  \item Manage organisation membership: generate invite codes, send email
        invites, and remove members.
  \item Assign and change member roles, and act on the intelligent
        role-recommendation banner on the Members settings page --- accepting
        a recommendation (which applies the role change under an
        optimistic-concurrency guard) or dismissing it (which starts a 30-day
        per-rule cool-down).
  \item Edit organisation settings (name, profile, and workspace-level
        configuration).
  \item Manage billing through Stripe~\cite{stripe}: view usage, manage
        payment methods, and review invoices produced by the usage-ticker and
        invoice-roller cron jobs; set the organisation's daily token budget on
        the cost tracker.
  \item View and export the append-only, hash-chained audit log (CSV export)
        and invoke its cryptographic verification endpoint to prove the chain
        has not been tampered with.
  \item Decide pending approvals from the recovery pipeline: \textbf{Approve}
        a validated patch, \textbf{Reject} it, or \textbf{Modify} it --- edit
        the unified diff in place before it ships, with the edited diff
        persisted as an RLHF feedback example.
  \item Configure each project's recovery policy, including
        rollback-on-SLO-breach behaviour.
  \item Operate the Docker preview-deploy engine from a project's
        \emph{Ops} tab: trigger Deploy and Stop, watch live status, open the
        returned preview URL, inspect build and container logs, and review
        deployment history.
  \item Install and manage integrations: the GitHub App (installation, repo
        listing, webhook-driven incident creation on merged pull requests),
        Sentry~\cite{sentry}, Datadog~\cite{datadog}, and
        PagerDuty~\cite{pagerduty} webhook sources, Slack~\cite{slackapi}
        (OAuth install plus interactive one-click approve/reject directly
        from a Slack message), and Argo CD~\cite{argocd}.
  \item Export evaluation-harness results as CSV (an admin-only action).
  \item Receive the daily admin digest e-mail summarising the last 24 hours
        of incidents, repairs, approvals, and MTTR for the organisation.
  \item Monitor platform operations pages: System Health (parallel dependency
        probes of PostgreSQL, Redis, Neo4j, MinIO, and Temporal), webhook
        activity, and integration connection health.
\end{itemize}

\subsection{Engineer/Member can:}

\begin{itemize}[nosep]
  \item Do everything a Viewer can do (role additivity).
  \item Sign up and sign in with e-mail and password (JWT session in an
        HttpOnly cookie), sign in passwordlessly via magic link, or sign in
        through WorkOS SSO~\cite{workos}; recover access through the
        password-reset flow (hashed, time-limited reset tokens).
  \item Enrol multi-factor authentication (TOTP with QR provisioning).
  \item Manage their own active sessions: list signed-in devices with IP
        address, user agent, and last-seen time (the caller's own session
        carries a ``This device'' badge), and revoke any session remotely.
  \item Create, list, and revoke scoped API keys (\code{nx\_live\_\ldots}
        bearer tokens) for programmatic access.
  \item Create projects through the new-project wizard and connect them to
        GitHub repositories.
  \item Browse the incidents list and open an incident's detail page, where
        the nine-agent pipeline timeline streams live over Server-Sent
        Events, including per-agent status, token counts, and cost.
  \item Inspect each agent's structured payloads (for example the Backend
        agent's raw generated diff) from the timeline.
  \item Trigger fault-injection runs from the Live-Demo console page, which
        presents 27 curated scenario cards drawn from a broader
        fixture-scenario taxonomy (one scenario is fully wired end-to-end
        today; the remainder are labelled ``Coming soon'').
  \item Run the evaluation harness: side-by-side OpenAI-versus-Ollama
        provider comparisons over fixture incidents, with persisted
        transcripts and cost figures.
  \item Query the knowledge base (pgvector-backed retrieval~\cite{pgvector})
        and view its retrieval status.
  \item Complete the post-incident NASA-TLX workload survey~\cite{nasatlx},
        the platform's instrument for measuring operator cognitive load.
  \item Use the operational consoles: the agents fleet page and per-agent
        detail views, the workflows list, the activity stream, the validator
        sandbox page, and the recovery-pipeline view.
\end{itemize}

\subsection{Viewer can:}

\begin{itemize}[nosep]
  \item Sign in and view the dashboard: KPI tiles (MTTR, recovery success
        rate, open incidents, mean tokens per run), the incidents-over-time
        chart, the system status panel, the projects health grid, and the
        recent activity feed.
  \item Read the incidents list and incident detail timelines, including the
        live SSE stream, without any decision or mutation controls.
  \item Read project pages (overview, activity, and status) and the
        performance page.
  \item View the agents, workflows, and activity-stream pages in read-only
        form.
  \item Take no mutating action of any kind: no approvals, no deployments, no
        settings changes, no key or session management beyond their own
        account.
\end{itemize}

\subsection{The 9-Agent Fleet autonomously:}

\begin{itemize}[nosep]
  \item \textbf{Detects} faults (Sentinel): an always-on detector applies two
        rules --- a per-row fatal-level trigger, and a statistical-process-%
        control spike rule that maintains a per-organisation EWMA mean and
        variance baseline and fires on a breach of mean~$+3\sigma$, with a
        Poisson floor for low-count noise and an absolute-floor warm-up
        fallback.
  \item \textbf{Diagnoses} root causes (Pathfinder): walks the Neo4j code
        dependency graph~\cite{neo4j} outward from the fault symptom and
        forwards candidate root-cause nodes, annotated with hop distance and
        degree, to the causal-inference sidecar, which ranks candidates by
        evidence-chain specificity, structural proximity, textual overlap,
        and a scenario prior, returning a per-component confidence
        breakdown~\cite{dowhy}.
  \item \textbf{Classifies and routes} (Synthesiser): assigns the incident a
        scenario (regex fast path with LLM fallback) and computes an ordered
        \code{selected\_agents} list; the workflow enforces this routing,
        skipping non-selected execution agents with a visible ``skipped''
        timeline state.
  \item \textbf{Plans} the repair (Architect): decomposes the incident into
        \code{plan\_steps}, \code{affected\_files}, and a
        \code{risk\_level}, and serves as a contract other agents may not
        silently violate --- if the generated diff touches a file outside
        \code{affected\_files}, the workflow flags a contract violation and
        forces HIGH severity so the patch can never auto-deploy unreviewed.
  \item \textbf{Codes} the fix (Backend): generates an LLM-guided minimal
        unified diff --- the smallest valid change rather than a full-file
        rewrite.
  \item \textbf{Tests} the fix (QA): auto-generates property-based test
        suites (pytest/Hypothesis~\cite{hypothesis}) from the incident and
        the diff; a background continuous-QA loop replays previously
        generated suites against the sandbox and raises a fresh incident on
        any pass-to-fail regression, closing a second, independent loop.
  \item \textbf{Packages} the deployment (DevOps): emits CI/CD and Argo CD
        manifests~\cite{argocd} and drives post-deployment health
        monitoring.
  \item \textbf{Handles data changes} (Data Engineer): proposes paired
        forward and reverse SQL migrations, runs a proactive schema-drift
        detector that introspects \code{information\_schema} against a
        captured baseline and raises a real incident on drift, and emits
        OpenLineage-compatible run events~\cite{openlineage} for every
        migration proposal and application.
  \item \textbf{Validates} every candidate patch (Validator): executes it in
        a locked-down Docker sandbox~\cite{docker} (no network, read-only
        filesystem, all capabilities dropped, PID-limited,
        no-new-privileges), runs the QA-generated property tests, and reports
        pass/fail and coverage; no unvalidated patch ever reaches the
        approval gate.
  \item \textbf{Classifies severity} at the approval gate: HIGH for
        schema-drift scenarios, sensitive-path matches (SQL, migrations,
        auth, security, crypto), unknown scenarios, or contract violations;
        MEDIUM otherwise, with a 120-second auto-approve countdown unless a
        human intervenes; LOW (tiny UI-only diffs) auto-approves immediately.
  \item \textbf{Deploys} approved patches (GitOps): opens a real GitHub pull
        request through a GitHub App installation
        (blob$\rightarrow$tree$\rightarrow$commit$\rightarrow$ref$\rightarrow$PR);
        when the triggering incident originated from the preview-deploy
        engine, a successful merge additionally triggers an automatic
        redeploy.
  \item \textbf{Rolls back} on regression: a post-deploy probe watches for a
        fresh fatal incident on the same organisation and service within a
        bounded window and, when the project policy permits, invokes Argo
        CD's rollback client to restore the prior known-good revision.
  \item \textbf{Records} every step as a Temporal-durable activity
        event~\cite{temporal} and streams it live to the incident detail
        page over Server-Sent Events.
  \item \textbf{Feeds the deploy engine's failures back into itself}: when a
        preview deployment fails to build or pass its health check, the
        deploy engine files a real incident (\code{source:
        "deploy\_engine"}) into the same Sentinel-to-pipeline path as every
        other incident source, making a broken build a self-healing target
        like any production fault.
\end{itemize}

\section{Non-Functional Requirements}

\begin{table}[htbp]
\caption{Summary of non-functional requirements and their realisation.}
\label{tab:nfr}
\centering
\begin{tabular}{p{3.4cm}p{9.6cm}}
\toprule
\textbf{Requirement} & \textbf{Realisation in Nexis} \\
\midrule
Multi-tenancy & PostgreSQL Row-Level Security on essentially every table~\cite{pgrls} \\
Durability & Temporal-orchestrated workflows survive crashes and restarts~\cite{temporal} \\
Security & Sandboxed patch execution, encrypted patch storage, hashed tokens, hash-chained audit log \\
Scalability & Stateless Go services, queue-decoupled workers, parallel dependency probes \\
Responsiveness & Server-Sent Events stream the live incident timeline \\
Usability & Role-scoped console, one-click approvals, plain-language rationales \\
\bottomrule
\end{tabular}
\end{table}

\textbf{Multi-tenancy and isolation.} Tenant isolation must not depend on
application code remembering to add a \code{WHERE org\_id = \ldots} clause.
Nexis enforces isolation in the database itself using PostgreSQL Row-Level
Security policies~\cite{pgrls}: a request-scoped transaction binds
\code{app.current\_org\_id} for every request, and the database refuses to
return or mutate rows belonging to any other organisation, regardless of what
the application layer asks for. This applies uniformly to human requests and
to agent writes.

\textbf{Durability of long-running work.} A recovery run spans many minutes,
crosses several services, waits on a human decision, and must survive process
crashes and restarts without losing its place. All pipeline state therefore
lives in Temporal~\cite{temporal}, a durable workflow engine: every agent
step is a journaled activity, and a resumed worker replays deterministically
to exactly where it stopped. An in-memory orchestrator was rejected because a
single crash mid-repair would strand an incident in an unknown state.

\textbf{Security.} Autonomously generated code is treated as untrusted by
construction. Candidate patches execute only inside a Docker
sandbox~\cite{docker} with networking disabled, a read-only filesystem, all
Linux capabilities dropped, a process-count limit, and
\code{no-new-privileges}. Patch blobs are stored encrypted in object storage.
Password-reset and API tokens are stored only as hashes (reset tokens are
32-byte random values, SHA-256 hashed, with a 30-minute TTL and a non-leaky
202 response). The audit log is append-only and hash-chained, so any
tampering with a past entry breaks the chain and is detectable through the
verification endpoint.

\textbf{Scalability.} The control plane, validator, GitOps, and deploy-engine
services are independent, stateless Go modules that scale horizontally;
Temporal decouples workflow load from any single process; Redis provides
rate limiting and caching; and the System Health page's dependency probes
fan out in parallel under a bounded timeout so a slow dependency cannot
stall the page.

\textbf{Real-time responsiveness.} An engineer watching a live repair must
see agent transitions as they happen, not on a polling delay. The incident
detail page subscribes to a Server-Sent Events stream, and every
Temporal-recorded activity event is pushed to the browser within moments of
occurring --- the live-demo runs in Chapter~6 were observed this way.

\textbf{Usability.} The console must make a genuinely complex system
operable: role-scoped navigation shows each user only what they can act on,
approvals are a one-click decision with the full diff in view, severity and
countdown state are always visible, and role recommendations carry a
plain-language rationale naming the concrete evidence behind them.

\section{System Feasibility}

\subsection{Technical Feasibility}

Nexis is deliberately polyglot, with each language chosen for the workload it
carries rather than for uniformity. The four backend services
(control-plane, validator, gitops, deploy-engine) are independent Go modules:
Go's static binaries, low-overhead concurrency, and first-class Docker API
client suit long-lived services that juggle webhooks, container lifecycles,
and parallel health probes. The console is Next.js/React with TypeScript,
which is the practical choice for a large role-gated UI with live SSE
streams, embedded Monaco diff editing, and chart-heavy dashboards. The
causal-inference sidecar is Python/FastAPI because the scientific-Python
ecosystem (and the Hypothesis property-testing stack~\cite{hypothesis} the QA
agent depends on) lives there.

The data layer pairs PostgreSQL 17 with pgvector~\cite{pgvector} --- one
engine provides transactional tenant data, row-level
security~\cite{pgrls}, and vector retrieval for the knowledge base, avoiding
a separate vector database --- while Neo4j~\cite{neo4j} holds the code
dependency graph, because root-cause diagnosis is a graph-traversal problem
(hop distances, degrees, paths from symptom to cause) that a relational
schema expresses poorly. Temporal~\cite{temporal} was proven in practice:
both live recovery runs reported in Chapter~6 executed durably end-to-end.
Every component is technology the team could run locally; the entire
platform was built and verified on a single development machine, with
Podman~\cite{podman} standing in for Docker~\cite{docker} as an
API-compatible runtime where Docker itself was unavailable.

\subsection{Economic Feasibility}

The platform's dependency set is open source end to end --- PostgreSQL,
Neo4j Community, Redis, MinIO, Temporal, OpenTelemetry~\cite{otel}, and the
Grafana observability stack --- so no licence cost attaches to development or
evaluation. The one economically significant input, LLM inference, is
dual-path by design: production targets OpenAI's hosted models, while
Ollama~\cite{ollama} serves fully supported local models
(\code{gpt-oss:20b}, \code{qwen2.5-coder:14b}) as a zero-marginal-cost
development and evaluation path. This is not a theoretical fallback: every
live end-to-end run in this report was driven by the local Ollama provider,
meaning the complete system --- including its most expensive component ---
was exercised without incurring any per-token cost. The cost tracker,
per-organisation token ledger, and daily budget controls exist precisely so
that the paid path remains bounded and observable when it is used.

\subsection{Operational Feasibility}

Operationally, the system is designed to be run by the people it serves. The
console presents role-based dashboards: an admin lands on organisation-wide
KPIs, approvals, and health; an engineer lands on incidents and the live
pipeline; a viewer sees a read-only picture of the same truth. High-stakes
autonomy is reduced to a one-click human decision (Approve/Reject/Modify),
optionally taken directly from Slack~\cite{slackapi} without opening the
console at all. Severity classification means routine low-risk fixes demand
no attention while sensitive changes always do. The NASA-TLX
instrument~\cite{nasatlx} is built in so that the operator-workload claim can
eventually be measured rather than asserted. Together these mean adopting
Nexis does not require a dedicated operations team; the intended day-to-day
operator is an ordinary engineer with a browser.
